% === Accuracy Without Attn2: Dependence on Current Value and Gap ===
% Appendix section
% Insert after Residual Stream Cleanup appendix

\section{Accuracy Without Attn2: Dependence on Current Value and Gap}
\label{sec:no-attn2-by-gap}

We ablate attn2 (skip the second-layer attention, retain MLP2) and measure
per-token accuracy as a function of the current sorted value~$i$ and the
gap~$g = v_{\mathrm{next}} - i$.  Figure~\ref{fig:no-attn2-by-gap} shows
three representative gaps: $g = 1$, $g = 10$, and $g = 50$.

For the smallest gap ($g = 1$), the model without attn2 is nearly at chance
when~$i$ is small, but accuracy rises steeply around $i \approx 150$ and
saturates close to~$1$ for $i \gtrsim 250$.  At $g = 10$ a similar ramp
appears but saturates at a lower level (${\approx}0.7$), and at $g = 50$
accuracy remains near zero for all values of~$i$.

This demonstrates that the first-layer pathway alone can resolve the correct
next token only when the target lies close in value to the current
token---a regime where the token embedding alone provides
enough information without the fine-grained windowed search that attn2
performs.  As the gap grows, attn2 becomes indispensable.

\begin{figure}[h]
  \centering
  \includegraphics[width=\linewidth]{newpics/no_attn2_acc_by_gap.png}
  \caption{Per-token accuracy without attn2 as a function of current sorted
    value~$i$, for gaps $g \in \{1, 10, 50\}$ (left to right).  Grey dots
    show per-$i$ accuracy; the solid curve is a rolling average (window 20);
    shaded bands show 95\% bootstrap confidence intervals.  For small gaps
    the first-layer pathway alone achieves high accuracy once~$i$ is
    sufficiently large; for large gaps attn2 is indispensable.}
  \label{fig:no-attn2-by-gap}
\end{figure}
